\documentclass[journal]{IEEEtran}
\usepackage[a5paper, margin=10mm, onecolumn]{geometry}
\usepackage{tfrupee}

\setlength{\headheight}{1cm}
\setlength{\headsep}{0mm}

\usepackage{gvv-book}
\usepackage{gvv}
\usepackage{cite}
\usepackage{amsmath,amssymb,amsfonts,amsthm}
\usepackage{algorithmic}
\usepackage{graphicx}
\usepackage{textcomp}
\usepackage{xcolor}
\usepackage{txfonts}
\usepackage{enumitem}
\usepackage{mathtools}
\usepackage{gensymb}
\usepackage{comment}
\usepackage[breaklinks=true]{hyperref}
\usepackage{tkz-euclide}

\graphicspath{{./figs/}}

\begin{document}
\title{4.13.79}
\author{AI25BTECH11010 - Dhanush Kumar}
\maketitle
\renewcommand{\thefigure}{\theenumi}
\renewcommand{\thetable}{\theenumi}

\textbf{Question}\\

In $\mathbb{R}^3$, let $L$ be a straight line passing through the origin. Suppose that all the points on $L$ are at a constant distance from two planes:$
P_1 : x+2y-z+1=0, 
P_2 : 2x-y+z-1=0.$
Let $M$ be the locus of the feet of the perpendicular drawn from the points on $L$ to the plane $P_1$. Which of the following points lie(s) on $M$?

\begin{enumerate}
\item $\myvec{0\\ -\tfrac{5}{6}\\ -\tfrac{2}{3}}$
\item $\myvec{-\tfrac{1}{6}\\ -\tfrac{1}{3}\\ \tfrac{1}{6}}$
\item $\myvec{-\tfrac{5}{6}\\ 0\\ \tfrac{2}{3}}$
\item $\myvec{-\tfrac{1}{3}\\ 0\\ \tfrac{2}{3}}$
\end{enumerate}

\textbf{Solution}\\

Step 1: Direction vector of $L$
\begin{align}
\vec{r}_L(t) &= t \vec{D}, \\
\vec{n}_1 &= \myvec{1\\2\\-1}, & \vec{n}_2 &= \myvec{2\\-1\\1}, \\
\vec{n}_1^T \vec{D} &= 0, & \vec{n}_2^T \vec{D} &= 0.
\end{align}

\begin{align}
	\myvec{\vec{n}_1 && \vec{n}_2}^T \vec{D} &= 0, \\
\vec{D} &= \myvec{1\\-3\\-5}, \\
\vec{r}_L(t) &= t \vec{D}.
\end{align}

Step 2: Foot of perpendicular to $P_1$
\begin{align}
\vec{r}_M &= \vec{r}_L(t) + \lambda \vec{n}_1, \\
&= t \vec{D} + \lambda \myvec{1\\2\\-1}, \\
\vec{n}_1^T \vec{r}_M + 1 &= 0, \\
\vec{n}_1^T \bigl( t \vec{D} + \lambda \myvec{1\\2\\-1} \bigr) + 1 &= 0, \\
\vec{n}_1^T \vec{r}_L(t) &= 1\cdot t + 2(-3t) + (-1)(-5t) = 0, \\
\vec{n}_1^T \vec{n}_1 &= 1^2 + 2^2 + (-1)^2 = 6, \\
6 \lambda + 1 &= 0 \;\Rightarrow\; \lambda = -\tfrac{1}{6}, \\
\vec{r}_M &= t \vec{D} - \tfrac{1}{6} \myvec{1\\2\\-1}, \\
\vec{r}_M &= \myvec{t-\tfrac{1}{6}\\ -3t-\tfrac{1}{3}\\ -5t+\tfrac{1}{6}}.
\end{align}

Step 3: Checking the options
\begin{align}
\text{Option (a):} \quad 
\myvec{0\\ -\tfrac{5}{6}\\ -\tfrac{2}{3}} &= \myvec{t-\tfrac{1}{6}\\ -3t-\tfrac{1}{3}\\ -5t+\tfrac{1}{6}} 
\;\Rightarrow\; t = \tfrac{1}{6}, \quad \text{lies on $M$},\\[2mm]
\text{Option (b):} \quad
\myvec{-\tfrac{1}{6}\\ -\tfrac{1}{3}\\ \tfrac{1}{6}} &= \myvec{t-\tfrac{1}{6}\\ -3t-\tfrac{1}{3}\\ -5t+\tfrac{1}{6}} 
\;\Rightarrow\; t = 0, \quad \text{lies on $M$},\\[2mm]
\text{Option (c):} \quad
\myvec{-\tfrac{5}{6}\\ 0\\ \tfrac{2}{3}} &= \myvec{t-\tfrac{1}{6}\\ -3t-\tfrac{1}{3}\\ -5t+\tfrac{1}{6}} 
\;\Rightarrow\; t=-\tfrac{2}{3}, \quad y = \tfrac{5}{3} \neq 0, \quad \text{does not lie on $M$},\\[2mm]
\text{Option (d):} \quad
\myvec{-\tfrac{1}{3}\\ 0\\ \tfrac{2}{3}} &= \myvec{t-\tfrac{1}{6}\\ -3t-\tfrac{1}{3}\\ -5t+\tfrac{1}{6}} 
\;\Rightarrow\; t=-\tfrac{1}{6}, \quad y = \tfrac{1}{6} \neq 0, \quad \text{does not lie on $M$}.
\end{align}

Final Answer
\begin{align}
\vec{r}_M &= \myvec{t-\tfrac{1}{6}\\ -3t-\tfrac{1}{3}\\ -5t+\tfrac{1}{6}}, \quad t \in \mathbb{R}, \\
\text{Points on $M$} &= \myvec{0\\ -\tfrac{5}{6}\\ -\tfrac{2}{3}}, \;\; \myvec{-\tfrac{1}{6}\\ -\tfrac{1}{3}\\ \tfrac{1}{6}}.
\end{align}

\end{document}
